\documentclass[12pt,a4paper]{article}
\usepackage[utf8]{inputenc}
\usepackage[spanish,es-tabla]{babel}
\usepackage{amsmath,amsfonts,amssymb}
\usepackage{graphicx}
\usepackage[left=3cm,right=3cm,top=3cm,bottom=3cm]{geometry}
\usepackage{titlesec}
\usepackage{hyperref}

% Configuración de hipervínculos (índice clickable)
\hypersetup{
    colorlinks=true,
    linkcolor=black,
    filecolor=magenta,      
    urlcolor=blue,
    pdftitle={Informe 01 de Proyecto de Tesis},
    pdfpagemode=FullScreen,
}

% Formato de los títulos (Secciones en Romanos, subsecciones en arábigos)
\titleformat{\section}{\large\bfseries}{\Roman{section}.}{1em}{}
\titleformat{\subsection}{\normalsize\bfseries}{\arabic{subsection}.}{1em}{}

\begin{document}

% ==============================================
% PORTADA
% ==============================================
\begin{titlepage}
    \centering
    
    {\Large \textbf{UNIVERSIDAD NACIONAL JORGE BASADRE GROHMANN}} \\[0.5cm]
    {\large FACULTAD DE CIENCIAS} \\[0.5cm] 
    {\large ESCUELA PROFESIONAL DE MATEMÁTICA} \\[2.5cm]
    
    {\Large \textbf{INFORME 01 DE PROYECTO DE TESIS}} \\[2cm]
    
    {\Large \textbf{Tema de Tesis:}} \\[0.5cm]
    {\large \textbf{Formulación y análisis de un modelo de programación no lineal para la optimización de costos en redes logísticas interregionales}} \\[3cm]
    
    {\large \textbf{Presentado por:}} \\[0.5cm]
    {\large Miguel Angel Luna Ttacca} \\[0.5cm]
    
    {\large \textbf{Asignatura:}} \\[0.5cm]
    {\large Seminario de Tesis I} \\[0.5cm]
    
    {\large \textbf{Docente:}} \\[0.5cm]
    {\large Dr. Rolando Wilman Vásquez Jaico} \\[1cm]
    
    {\large Tacna - Perú} \\[0.5cm]
    {\large 2026}
    
\end{titlepage}

% ==============================================
% ÍNDICE GENERAL
% ==============================================
\tableofcontents
\newpage

% ==============================================
% DESARROLLO DEL INFORME
% ==============================================

\section{Preparación y revisión bibliográfica y definición del problema}

\subsection{Búsqueda y selección de literatura relevante}
Para la construcción del marco teórico y el estado del arte de la presente investigación, se realizó una búsqueda exhaustiva dividida en dos ejes fundamentales. El primer eje corresponde a los fundamentos matemáticos de la optimización, para lo cual se seleccionaron textos clásicos de rigor axiomático, tales como \textit{Nonlinear Programming: Theory and Algorithms} (Bazaraa, Sherali \& Shetty, 2006), \textit{Numerical Optimization} (Nocedal \& Wright, 2006) y \textit{Convex Optimization} (Boyd \& Vandenberghe, 2004). El segundo eje se centró en la aplicación de modelos a redes de transporte, recurriendo a bases de datos académicas (Redalyc, Dialnet, SciELO y repositorios universitarios), seleccionando artículos que abordan la optimización de procesos de transporte de carga terrestre e integración de producción en redes logísticas.

\subsection{Lectura y análisis de la literatura}
El análisis de la literatura revela que el estudio de redes logísticas se apoya históricamente en la teoría de flujos en redes (Ahuja, Magnanti \& Orlin, 1993). En el ámbito aplicado, estudios recientes demuestran el esfuerzo continuo por modelar el transporte primario y la integración de inventarios (como los casos estudiados en el sector alimentos en Colombia y modelos de integración plantas-distribuidores). Sin embargo, desde la perspectiva matemática pura, gran parte de esta literatura aplicada recurre a la Programación Lineal (PL) o a la Programación Lineal Entera Mixta (PLEM) asumiendo proporcionalidad estricta en los costos, o bien, emplean algoritmos heurísticos que no garantizan la convergencia hacia un óptimo global y cuyo análisis de complejidad suele ser omitido.

\subsection{Identificación de brechas y problemas de investigación}
Se ha identificado una brecha teórica significativa: mientras que la ingeniería logística busca soluciones rápidas asumiendo linealidad, la realidad física del transporte interregional involucra dinámicas intrínsecamente no lineales (ej. economías de escala decrecientes, costos exponenciales por congestión de rutas y tarifas de almacenamiento por tramos). La literatura aplicada revisada carece, en su mayoría, de un análisis riguroso sobre la topología del espacio de soluciones cuando se introducen estas restricciones no lineales. Falta un puente formal que no solo plantee la ecuación empírica, sino que demuestre axiomáticamente la existencia de soluciones, analice la convexidad del modelo y justifique numéricamente el algoritmo resolutivo empleado.

\subsection{Definición y formulación del problema de investigación en términos matemáticos}
Sea un grafo dirigido $G = (V, E)$, donde $V$ es el conjunto de nodos (centros de producción, transbordo y demanda) y $E$ es el conjunto de arcos (rutas interregionales). Sea $x_{ij} \in \mathbb{R}^{+}$ la variable de decisión continua que representa el flujo de carga a través del arco $(i,j) \in E$. 

El problema consiste en minimizar una función objetivo multivariable no lineal $F(X)$:
\begin{equation}
    \min_{X \in \Omega} F(X) = \sum_{(i,j) \in E} f_{ij}(x_{ij}) + \sum_{i \in V} h_i(y_i)
\end{equation}
donde $f_{ij}$ es una función de costo de transporte estrictamente no lineal (posiblemente no convexa debido a economías de escala), $h_i$ representa el costo de almacenamiento no lineal en el nodo $i$, y $y_i$ es el inventario en dicho nodo.

El problema está sujeto a la región factible $\Omega \subset \mathbb{R}^{|E|}$, definida por el sistema de restricciones:
\begin{equation}
    \sum_{j: (i,j) \in E} x_{ij} - \sum_{j: (j,i) \in E} x_{ji} = d_i, \quad \forall i \in V \quad \text{(Conservación de flujo)}
\end{equation}
\begin{equation}
    0 \leq x_{ij} \leq u_{ij}, \quad \forall (i,j) \in E \quad \text{(Capacidad de rutas)}
\end{equation}
\begin{equation}
    g_k(X) \leq 0, \quad \forall k \in K \quad \text{(Restricciones interregionales complejas)}
\end{equation}

El desafío matemático radica en que, dado que $F(X)$ y $g_k(X)$ no son lineales, la región $\Omega$ puede perder sus propiedades de poliedro convexo, lo que dificulta establecer las condiciones necesarias y suficientes de optimalidad global mediante los teoremas clásicos.

\subsection{Planteamiento de las preguntas de investigación que se buscan responder}
\textbf{Pregunta principal:}
¿De qué manera la formulación y el análisis de un modelo de programación no lineal permite la optimización de costos en redes logísticas interregionales?

\textbf{Preguntas secundarias:}
\begin{itemize}
    \item ¿Cómo se formulan analítica y axiomáticamente las funciones multivariables y el sistema de restricciones para representar los costos operativos no lineales en un grafo dirigido $G=(V,E)$?
    \item ¿Qué propiedades geométricas y topológicas, y bajo qué condiciones de optimalidad (KKT), rigen la región factible del modelo propuesto para garantizar la identificación de óptimos locales o globales?
    \item ¿Cuál es la tasa de convergencia asintótica y el grado de complejidad computacional de los métodos numéricos de optimización continua aplicados a la resolución de este modelo matemático?
\end{itemize}

\subsection{Explicación de los problemas secundarios que se abordarán}
Para dar respuesta a las preguntas planteadas, la investigación abordará teóricamente los siguientes problemas secundarios, derivados directamente de las preguntas secundarias:
\begin{enumerate}
    \item \textbf{Problema de Formulación Analítica y Axiomática:} Se abordará la dificultad de representar matemáticamente las no linealidades de los costos operativos (como economías de escala y congestión) mediante funciones multivariables continuas y sistemas de restricciones en un grafo dirigido $G=(V,E)$, garantizando el cumplimiento de los principios de conservación de flujo.
    \item \textbf{Problema de Análisis Geométrico, Topológico y Optimalidad:} Se enfrentará el desafío de determinar las propiedades de compacidad y convexidad de la región factible, así como derivar y analizar las condiciones de Karush-Kuhn-Tucker (KKT) necesarias para distinguir y garantizar la existencia de óptimos locales o globales en el modelo.
    \item \textbf{Problema de Resolución Numérica y Complejidad Computacional:} Se tratará el reto de seleccionar, aplicar y evaluar métodos numéricos de optimización continua, determinando analíticamente su tasa de convergencia asintótica y el grado de complejidad computacional requerido para resolver el modelo de manera eficiente.
\end{enumerate}

\section{Antecedentes del problema}

\subsection{Revisión de trabajos previos relacionados con el problema}
El estudio de la optimización en redes logísticas ha sido un foco recurrente tanto en la investigación de operaciones como en la matemática aplicada. En el ámbito regional latinoamericano, investigaciones recientes han modelado redes de transporte primario; por ejemplo, estudios enfocados en el sector de alimentos en Colombia han analizado cómo la estructuración de un modelo matemático estático minimiza los costos operativos. Asimismo, existen antecedentes que proponen modelos matemáticos de integración de producción e inventarios (conectando plantas, distribuidores y detallistas), los cuales abordan la red como un sistema multi-escalón (\textit{multi-echelon}) con restricciones acopladas de almacenamiento.

En la literatura referente al diseño de redes logísticas inversas y la optimización de procesos de transporte de carga terrestre, la constante metodológica ha sido la formulación de problemas de Programación Lineal (PL) o Programación Lineal Entera Mixta (PLEM). Estos trabajos asumen, por conveniencia de la resolución computacional, que las funciones de costo crecen de forma estrictamente lineal o escalonada. Aunque estos estudios arrojan resultados prácticos útiles, desde una perspectiva matemática rigurosa, incurren en una pérdida significativa de precisión al ignorar el comportamiento inherentemente no lineal de dinámicas reales, como la congestión geométrica de rutas vehiculares y las curvas cóncavas de las economías de escala.

\subsection{Evolución histórica del problema y su estado actual en la literatura}
Históricamente, el modelamiento algebraico de redes de transporte tiene sus raíces formales en el problema clásico de transporte de Kantorovich y Hitchcock (década de 1940). En las décadas de 1960 y 1970, la maduración de la teoría algorítmica de grafos permitió el desarrollo de modelos de flujo de costo mínimo con bases matemáticas sólidas. 

Durante los años 1990 y principios de los 2000, impulsados por el aumento del poder computacional, los investigadores transicionaron masivamente hacia modelos discretos y combinatorios (PLEM) para representar variables binarias topológicas (como la apertura de arcos). Sin embargo, el \textbf{estado actual de la literatura} en matemática de frontera reconoce que las macrorredes modernas se comportan como flujos fluidodinámicos continuos. Por tanto, el paradigma investigativo reciente está rotando hacia la \textbf{Programación No Lineal Continua}, que permite incorporar funciones diferenciables no convexas para representar variaciones tarifarias marginales, abriendo nuevos problemas teóricos respecto a la existencia de múltiples óptimos locales.

\subsection{Principales teorías, métodos y resultados previos que se han obtenido}
Los modelos desarrollados hasta la fecha se fundamentan en un andamiaje teórico que combina la Teoría Algebraica de Grafos (mediante el uso de matrices de incidencia nodo-arco para las restricciones topológicas) y el Análisis Convexo (para la caracterización de la región factible).

En cuanto a los \textbf{métodos resolutivos} empleados en los antecedentes, se observan tres tendencias claras:
\begin{itemize}
    \item \textbf{Métodos Lineales Exactos:} Utilización de Descomposición de Benders y métodos Branch-and-Bound, los cuales padecen de explosión combinatoria (complejidad NP-Hard) ante el aumento de la densidad del grafo.
    \item \textbf{Metaheurísticas de Caja Negra:} Implementadas ampliamente en artículos orientados a la ingeniería pura. Si bien obtienen tiempos computacionales bajos, carecen de demostraciones formales de convergencia topológica hacia un óptimo global.
    \item \textbf{Métodos de Optimización Diferenciable (PNL):} Aplicación de métodos de descenso por gradiente y Programación Cuadrática Secuencial (SQP). La literatura matemática exige que las funciones estudiadas sean de clase $C^2$ (dos veces diferenciables) para asegurar el cálculo del Hessiano.
\end{itemize}

Los \textbf{resultados previos} evidencian que incorporar funciones de costo no lineales incrementa drásticamente la dificultad analítica y el tiempo de convergencia de los algoritmos; no obstante, minimiza sustancialmente la brecha de error entre las predicciones del modelo abstracto y las mediciones empíricas de los costos logísticos.

\section{Objetivos e hipótesis de la investigación}

\subsection{Declaración clara del objetivo principal de la investigación}
Formular y analizar rigurosamente un modelo matemático fundamentado en la Programación No Lineal para la minimización de costos en redes logísticas interregionales, estableciendo sus propiedades topológicas subyacentes y evaluando teóricamente la convergencia computacional de algoritmos de optimización diferenciable aplicados a dicho modelo.

\subsection{Desglose de objetivos específicos que se pretenden alcanzar para abordar el problema}
Para garantizar el cumplimiento del objetivo principal y asegurar el peso matemático de la investigación, se establecen los siguientes objetivos específicos:

\begin{enumerate}
    \item \textbf{Construcción Axiomática del Modelo:} Definir analíticamente las funciones multivariables continuas (función objetivo y sistema de restricciones) que representen las no linealidades de los costos operativos (economías de escala y congestión) proyectadas sobre un grafo dirigido $G = (V, E)$, respetando estrictamente los principios de conservación de flujo matricial.
    
    \item \textbf{Análisis Geométrico y Topológico:} Demostrar teóricamente las propiedades de compacidad y convexidad (o falta de ella) de la región factible generada, procediendo a derivar y analizar las condiciones necesarias y suficientes de optimalidad de Karush-Kuhn-Tucker (KKT) para la identificación de óptimos locales frente a óptimos globales.
    
    \item \textbf{Análisis Numérico y Computacional:} Implementar computacionalmente métodos numéricos de optimización continua de orden superior (como el Método de Puntos Interiores o algoritmos Cuasi-Newton) para la resolución del modelo, evaluando y demostrando matemáticamente su tasa de convergencia asintótica y su grado de complejidad computacional.
\end{enumerate}

\subsection{Hipótesis de trabajo}
Al tratarse de una investigación de corte teórico-práctico en el campo de la matemática aplicada, se postula la siguiente hipótesis fundamental que guía el desarrollo del análisis:

\textbf{Hipótesis Principal:}
La formulación de una red logística interregional utilizando funciones de costo de clase $C^2$ (dos veces diferenciables continuas) y no estrictamente convexas, configura un espacio topológico de soluciones donde la aplicación de algoritmos de optimización basados en el gradiente conjugado o Puntos Interiores logra converger hacia óptimos estacionarios con un menor error de aproximación analítica en comparación a los modelos de Programación Lineal Entera Mixta (PLEM), asumiendo un incremento en la complejidad temporal acotado por un límite polinomial estable.

\section{Justificación del Problema}

\subsection{Razones por las cuales el problema es relevante e importante}
En el contexto moderno del flujo masivo de mercancías, las redes logísticas interregionales han adquirido complejidades dinámicas sin precedentes. A pesar de esto, gran parte de las metodologías empleadas para su diseño se basan en simplificaciones lineales que asumen costos de transporte y almacenamiento directamente proporcionales. Esta linealización sistemática omite fenómenos inherentes a la realidad física, tales como el encarecimiento exponencial por la congestión vehicular en ciertas rutas (comportamiento convexo) y la reducción marginal de costos por volumen gracias a las economías de escala (comportamiento cóncavo). Abordar este problema mediante Programación No Lineal es crucialmente relevante porque cierra la brecha empírica entre el modelo abstracto y el comportamiento real del sistema, evitando óptimos "artificiales" generados por la excesiva simplificación del problema original.

\subsection{Impacto potencial de resolver el problema en el campo de la matemática teórica}
Desde la perspectiva de la matemática pura, y en particular dentro de la Optimización Continua y el Análisis Convexo, el estudio de flujos en redes gobernados por funciones no lineales aporta a la comprensión de la topología en espacios vectoriales restringidos. Al introducir curvas de costo no estrictamente convexas sobre un grafo, se desafían las asunciones clásicas de unicidad. Resolver matemáticamente este problema demanda demostrar formalmente bajo qué condiciones geométricas la región factible permite el cumplimiento robusto de las condiciones de Karush-Kuhn-Tucker (KKT). Asimismo, el análisis del problema dual asociado aportará conocimiento teórico sobre el comportamiento y la estabilidad de los multiplicadores de Lagrange (precios sombra) ante perturbaciones paramétricas en sistemas no lineales de gran escala.

\subsection{Beneficios teóricos y prácticos de abordar este problema}
La estructuración de este modelo genera aportes duales que se alinean perfectamente con la rama de la Matemática Pura y Aplicada:

\textbf{Beneficios Teóricos:}
\begin{itemize}
    \item \textbf{Sustento Analítico:} Provee un marco de demostraciones formales que justifican algebraicamente el uso de algoritmos de optimización basada en gradientes sobre grafos, contrastando con el uso empírico de metaheurísticas de "caja negra" comunes en ingeniería.
    \item \textbf{Cotas Computacionales:} Aporta al estudio del análisis numérico al evaluar y demostrar la tasa de convergencia asintótica y el comportamiento del error matemático de métodos como Puntos Interiores o Cuasi-Newton frente a matrices Jacobianas y Hessianas dispersas (típicas de las topologías de red).
\end{itemize}

\textbf{Beneficios Prácticos:}
\begin{itemize}
    \item \textbf{Resolución Computacional Aplicada:} Deriva en el desarrollo de un algoritmo numérico programable que podría ser implementado por tomadores de decisión para calcular ruteos y distribuciones volumétricas que minimicen genuinamente el costo económico absoluto.
    \item \textbf{Eficiencia Sistémica:} Al optimizar los flujos considerando restricciones no lineales (como la capacidad real elástica de las rutas), se reduce el riesgo de congestiones logísticas, lo que disminuye el consumo de combustible y, colateralmente, la huella de carbono de las flotas interregionales.
\end{itemize}

\section{Alcance y Limitaciones}

\subsection{Definición del alcance de la investigación}
La presente investigación tiene un alcance teórico-computacional enmarcado estrictamente en el dominio de la Optimización Continua y la Matemática Aplicada. A nivel teórico, el trabajo comprenderá la formulación algebraica del modelo matemático no lineal sobre redes, la demostración rigurosa de sus propiedades geométricas (convexidad del espacio funcional), y la estructuración de las condiciones analíticas de optimalidad de Karush-Kuhn-Tucker. A nivel aplicativo, el estudio abarcará la selección, codificación y evaluación asintótica de un algoritmo numérico de descenso iterativo, culminando en la simulación computacional de diversas instancias abstractas del problema logístico para medir experimentalmente la tasa de convergencia del método matemático propuesto.

\subsection{Identificación de las limitaciones y restricciones que puedan afectar el estudio}
La consecución rigurosa de los objetivos de la tesis podría verse condicionada por las siguientes limitaciones intrínsecas:
\begin{itemize}
    \item \textbf{Explosión Computacional (Complejidad):} Dado que la pérdida de la estricta convexidad en problemas multivariables eleva drásticamente el costo computacional por iteración, la validación asintótica del algoritmo estará limitada en escala. La cardinalidad de los conjuntos del grafo (nodos y arcos) que se puedan someter a simulación dependerá de las capacidades de la memoria RAM y el hardware de cómputo disponible.
    \item \textbf{Acceso a Parámetros Industriales Reales:} Extraer coeficientes reales para parametrizar las ecuaciones continuas $C^2$ (por ejemplo, el parámetro de decaimiento marginal en una economía de escala real) es complejo debido a acuerdos de confidencialidad en el sector transporte. Por ello, la validación del modelo podría requerir la generación algorítmica de instancias de prueba o la adaptación de \textit{benchmarks} académicos estandarizados provenientes de la literatura internacional.
\end{itemize}

\subsection{Delimitación de lo que se va a incluir y excluir en la investigación}
Para garantizar la viabilidad del estudio, resguardar el rigor de las demostraciones matemáticas y delimitar el área de experticia evaluada, se establecen las siguientes condiciones de frontera:

\textbf{El modelo INCLUIRÁ:}
\begin{itemize}
    \item Variables de decisión pertenecientes estrictamente al campo de los números reales no negativos ($\mathbb{R}^+$), enfocándose en modelos de variable continua.
    \item Funciones objetivo y de restricción algebraicas que cumplan con ser de clase $C^2$ (diferenciables hasta su segunda derivada) para permitir el cálculo formal del vector Gradiente y la matriz Hessiana.
    \item Un enfoque estático-determinista de red macroscópica (el volumen de la demanda en el nodo de destino es un escalar conocido a priori, y el arco representa una vía interregional).
\end{itemize}

\textbf{El modelo EXCLUIRÁ expresamente:}
\begin{itemize}
    \item \textbf{Incertidumbre Matemática (Programación Estocástica):} No se incorporarán variables aleatorias o funciones de densidad de probabilidad para modelar la fluctuación de la demanda en el tiempo. Todas las matrices de datos se asumirán invariables y deterministas.
    \item \textbf{Variables Discretas / Combinatoria (MINLP):} Se excluirá el modelado con variables enteras o binarias booleanas (ej. abrir o cerrar un centro). En caso de plantearse restricciones binarias por necesidad topológica, estas se abordarán mediante su respectiva \textit{relajación continua}.
    \item \textbf{Micro-ruteo Vehicular (VRP):} El modelo no abordará problemas como el del agente viajero (TSP) o el ruteo interno de camiones a nivel de calles urbanas, ya que el alcance topológico se ciñe exclusivamente a los grandes flujos entre nodos macro-regionales.
\end{itemize}

\section{Metodología Resumida}

\subsection{Breve descripción de los métodos y enfoques que se utilizarán para abordar el problema}
La investigación se estructura bajo un paradigma cuantitativo con un diseño metodológico integral que se bifurca en dos grandes fases: una analítico-deductiva (correspondiente al rigor de la Matemática Pura) y otra numérico-experimental (orientada a la Matemática Aplicada).

\textbf{Fase Analítico-Deductiva (Formulación Teórica):}
Se empleará rigurosamente el \textbf{método lógico-deductivo}. Se partirá de los axiomas de los espacios vectoriales reales y de la teoría algebraica de grafos para estructurar el problema. Los procedimientos analíticos a utilizar incluyen:
\begin{enumerate}
    \item \textbf{Cálculo Diferencial Multivariable:} Para el cálculo exacto del vector Gradiente $\nabla F(X)$ y la matriz Hessiana $\nabla^2 F(X)$ de las funciones de costo logístico.
    \item \textbf{Análisis Topológico y Convexo:} Empleo de teoremas clásicos (como el Teorema de Weierstrass) para evaluar la compacidad del espacio de soluciones $\Omega$ y demostrar formalmente la existencia del óptimo.
    \item \textbf{Derivación de Dualidad:} Formulación analítica de las condiciones necesarias y suficientes de optimalidad de Karush-Kuhn-Tucker (KKT).
\end{enumerate}

\textbf{Fase Numérico-Experimental (Resolución Computacional):}
En la etapa de aplicación, el método radicará en el \textbf{diseño y experimentación in-silico} algorítmica:
\begin{enumerate}
    \item \textbf{Implementación Numérica:} Codificación de métodos iterativos de optimización diferenciable (como el Método de Puntos Interiores o la Programación Cuadrática Secuencial - SQP) en lenguajes orientados al cálculo matricial de alto rendimiento (ej. Julia, Python/SciPy).
    \item \textbf{Diseño de Experimentos Algorítmicos:} Ejecución de simulaciones numéricas sobre una batería de grafos (instancias de prueba) generadas paramétricamente, con el objetivo de cuantificar variables de rendimiento asintótico (número de iteraciones para satisfacer la tolerancia del error residual $\epsilon$, y costo temporal de CPU).
\end{enumerate}

\subsection{Justificación de la elección de estos métodos}
La elección del \textbf{método deductivo} formal para la formulación teórica es indiscutible e irrenunciable dentro del ámbito académico de las Ciencias Matemáticas. Mientras que otras disciplinas toleran la validación empírica o heurística, en matemática pura la única forma de garantizar que una solución propuesta es correcta (y no es un simple óptimo local engañoso) es a través del blindaje de una demostración analítica irrefutable sobre la topología del problema.

Respecto a la fase aplicada, la selección de los \textbf{métodos iterativos continuos (Puntos Interiores y SQP)} en contraposición a las tradicionales "metaheurísticas de ingeniería" (como los Algoritmos Genéticos) obedece a una necesidad de rigor. Las metaheurísticas son cajas negras estocásticas que carecen de pruebas analíticas de convergencia. En contraste, los métodos matriciales iterativos seleccionados se fundamentan teóricamente en las aproximaciones del Teorema de Taylor y el sistema de KKT, lo que permite demostrar matemáticamente su tasa de convergencia (cuadrática o superlineal) y acotar la cota del error algebraico en cada paso del algoritmo.

\section{Relevancia del Estudio}

\subsection{Impacto potencial del estudio en la teoría matemática}
El impacto principal de esta investigación radica en su enriquecimiento teórico al campo de la Optimización Continua y el Análisis Convexo aplicado a grafos. Al trascender el paradigma de las aproximaciones lineales, el estudio proporciona demostraciones analíticas que clarifican el comportamiento geométrico de regiones factibles limitadas por inecuaciones no lineales. Específicamente, el escrutinio sobre la matriz Hessiana de las funciones de costo y el análisis de sensibilidad de los multiplicadores de Lagrange (en el marco de las condiciones de Karush-Kuhn-Tucker) ofrecerá nuevas perspectivas matemáticas sobre cómo las no linealidades originadas por economías de escala afectan la compacidad del espacio de soluciones y la unicidad de los óptimos globales.

\subsection{Contribución esperada del estudio a la resolución de problemas matemáticos más amplios}
El aparato deductivo y los resultados algorítmicos desarrollados en esta tesis trascenderán la particularidad del grafo logístico y servirán de cimiento para el abordaje de problemas matemáticos de escala macro. La evaluación teórica de la complejidad temporal y la cota del error numérico en algoritmos de segundo orden frente a matrices dispersas (como la matriz de incidencia del grafo) arroja luz sobre la eficiencia computacional de los Métodos de Puntos Interiores en sistemas multivariables. Estas demostraciones son altamente valoradas y extrapolables para la resolución de la clase abstracta de problemas de "Flujo de Múltiples Mercancías con Costo No Lineal" (\textit{Non-linear Multi-commodity Flow Problems}), cuyo estudio de intratabilidad (espectro NP-Hard) es un área hiperactiva en las matemáticas computacionales.

\subsection{Posibles aplicaciones de los resultados en otras áreas de la matemática o en campos relacionados}
Debido a que el modelo matemático se construye axiomáticamente sobre una estructura topológica abstracta $G=(V,E)$, las ecuaciones, teoremas y el algoritmo iterativo subyacente pueden ser despojados de su contexto logístico para solucionar problemas isomórficos en disciplinas afines:
\begin{itemize}
    \item \textbf{Teoría de Redes de Telecomunicaciones:} El ruteo óptimo de paquetes de datos a través de servidores de internet, donde la congestión del tráfico genera retardos exponenciales (costo no lineal), se rige por las mismas ecuaciones de conservación de flujo y capacidad diferencial modeladas en esta tesis.
    \item \textbf{Mecánica de Fluidos e Hidráulica de Redes:} El cálculo de caudales óptimos en sistemas de gasoductos o acueductos interregionales. En estos sistemas, el principio de continuidad (nodos) y las ecuaciones de pérdida de carga por fricción (arcos) constituyen intrínsecamente un problema no lineal equivalente.
    \item \textbf{Ciencia de Datos y Machine Learning:} El estudio algorítmico del descenso por gradientes en estructuras de red es el soporte matemático fundamental de las Redes Neuronales de Grafos (\textit{Graph Neural Networks - GNN}), permitiendo que las estrategias de minimización desarrolladas aquí sean adaptables a arquitecturas de aprendizaje profundo.
\end{itemize}

\section{Conclusiones}

\subsection{Resumen de la importancia del problema y de la necesidad de su estudio}
La presente investigación subraya la necesidad imperativa de trascender las simplificaciones lineales en el modelado de redes logísticas interregionales. La importancia del problema radica en que el comportamiento real de los costos de transporte y almacenamiento está intrínsecamente ligado a fenómenos no lineales, como las economías de escala y la saturación de infraestructura. Ignorar estas dinámicas no solo conduce a soluciones subóptimas, sino que genera una desconexión entre la teoría matemática y la realidad operativa. El estudio de este problema es esencial para proporcionar un marco analítico robusto que permita capturar la complejidad de los sistemas logísticos modernos, garantizando que la optimización se base en representaciones matemáticas fieles a la naturaleza física y económica del flujo de bienes.

\subsection{Conexión entre el planteamiento del problema y los objetivos de la investigación}
Existe una correspondencia directa y necesaria entre el planteamiento del problema y los objetivos definidos. Ante la brecha identificada entre los modelos lineales tradicionales y la realidad no lineal, la investigación se ha planteado objetivos que abordan esta problemática desde tres frentes complementarios:
\begin{enumerate}
    \item La construcción de un modelo algebraico riguroso responde directamente a la necesidad de formalizar las dinámicas de costo no lineales sobre una estructura de grafo.
    \item El análisis topológico y la derivación de las condiciones de Karush-Kuhn-Tucker (KKT) permiten fundamentar teóricamente la viabilidad de encontrar soluciones óptimas, proporcionando el rigor matemático que el planteamiento del problema exige.
    \item Finalmente, la implementación y evaluación de algoritmos numéricos de optimización continua cierran el ciclo de investigación, transformando el problema abstracto en una herramienta de resolución capaz de validar la eficiencia y convergencia del modelo propuesto. 
\end{enumerate}
De esta manera, el cumplimiento de los objetivos garantiza una respuesta integral a las preguntas de investigación planteadas al inicio del estudio.


\newpage
\addcontentsline{toc}{section}{Referencias Bibliográficas}
\begin{thebibliography}{99}

\bibitem{ahuja1993} 
Ahuja, R. K., Magnanti, T. L., \& Orlin, J. B. (1993). \textit{Network Flows: Theory, Algorithms, and Applications}. Prentice Hall.

\bibitem{bazaraa2006} 
Bazaraa, M. S., Sherali, H. D., \& Shetty, C. M. (2006). \textit{Nonlinear Programming: Theory and Algorithms} (3ra ed.). John Wiley \& Sons.

\bibitem{boyd2004} 
Boyd, S., \& Vandenberghe, L. (2004). \textit{Convex Optimization}. Cambridge University Press.

\bibitem{inversa} 
\textit{Diseño de redes de logística inversa: una revisión del estado del arte y aplicación práctica}. (s.f.). Recuperado de \url{https://scispace.com/pdf/diseno-de-redes-de-logistica-inversa-una-revision-del-estado-1xhlml728v.pdf}

\bibitem{transporte} 
\textit{Influencia de modelos matemáticos para la optimización de los procesos de transporte de carga terrestre}. (s.f.). Revista CYT. Recuperado de \url{https://revistas.unac.edu.pe/index.php/CYT/article/view/106}

\bibitem{colombia} 
\textit{Modelo de optimización de la red de transporte primario en el sector de alimentos en Colombia}. (s.f.). Redalyc. Recuperado de \url{https://www.redalyc.org/journal/2856/285667184007/html/}

\bibitem{inventarios} 
\textit{Modelo matemático de integración de producción e inventarios en una red logística plantas, distribuidores y detallistas}. (s.f.). Dialnet. Recuperado de \url{https://dialnet.unirioja.es/descarga/articulo/8704880.pdf}

\bibitem{nocedal2006} 
Nocedal, J., \& Wright, S. J. (2006). \textit{Numerical Optimization} (2da ed.). Springer.

\bibitem{ravindran2012} 
Ravindran, A. R., \& Warsing, D. P. (2012). \textit{Supply Chain Engineering: Models and Applications}. CRC Press.

\end{thebibliography}

\end{document}
